\section{Theoretical Analysis}
\label{sec:theory}

Let \(w^\star\in\arg\min_{w\in\Delta_K}R(w)\).  Our default guarantee needs
only a one-sided condition at this fixed comparator.

\begin{assumption}[Comparator margin condition]
\label{ass:comp-margin}
There are \(C_m>0\), \(\beta>0\), and \(t_0>0\) such that
\[
\mathbb P_{q\sim\mathcal D}
\{0<M_q(w^\star)\le t\}
\le C_m t^\beta,
\qquad 0<t\le t_0.
\]
\end{assumption}

This condition controls only correctly solved problems that are close to a
tie.  It says nothing about negative margins and, importantly, imposes no
margin condition at the random estimator.

\subsection{Uniform control of finite-generation margins}
\label{sec:theory-concentration}

\begin{lemma}[Uniform training-sample margin accuracy]
\label{lem:margin-accuracy}
For \(\delta\in(0,1)\), set
\(\varepsilon_N(\delta)=\sqrt{2\log(8KAQ/\delta)/N}\).  With probability at
least \(1-\delta/2\) over the generations, conditional on the sampled questions,
\[
\sup_{j\le Q}\sup_{w\in\Delta_K}
|\widehat M_{q_j}(w)-M_{q_j}(w)|
\le\varepsilon_N(\delta).
\]
\end{lemma}

\begin{proof}
Hoeffding's inequality and a union bound over at most \(KAQ\) answer
frequencies give
\[
\max_{j,i,a}
|\widehat p_i(a\mid q_j)-p_i(a\mid q_j)|
\le\frac{\varepsilon_N}{2}
\]
with probability at least \(1-\delta/2\).  Every coordinate of an answer-gap
vector then has error at most \(\varepsilon_N\).  Since
\(\|w\|_1=1\), each weighted gap has the same bound; stability of the minimum
over wrong answers preserves it.
\end{proof}

On this event, for every training question and every \(w\),
\begin{equation}
\mathbf 1\{M_{q_j}(w)\le0\}
\le
\mathbf 1\{\widehat M_{q_j}(w)\le\varepsilon_N\}
\le
\mathbf 1\{M_{q_j}(w)\le2\varepsilon_N\}.
\label{eq:indicator-sandwich}
\end{equation}
This uniform sandwich is the reason for introducing the buffer.

For \(t\ge0\), define
\(\mathcal F_t=\{q\mapsto\mathbf 1\{M_q(w)\le t\}:w\in\Delta_K\}\).
On \(Q\) fixed questions, at most \(Q(A-1)\) affine hyperplanes partition the
\((K-1)\)-dimensional simplex.  Thus
\[
\Pi_{\mathcal F_t}(Q)
\le
\sum_{r=0}^{K-1}\binom{Q(A-1)}r,
\]
and standard VC symmetrization yields, simultaneously for
\(t\in\{0,2\varepsilon_N\}\),
\begin{equation}
\sup_{w\in\Delta_K}
\left|
\frac1Q\sum_{j=1}^Q\mathbf 1\{M_{q_j}(w)\le t\}
-
\mathbb P\{M_q(w)\le t\}
\right|
\le
\mathfrak v_Q(\delta),
\label{eq:vq-event}
\end{equation}
where
\[
\mathfrak v_Q(\delta)
=
C\sqrt{
\frac{(K-1)\log(eAQ)+\log(1/\delta)}{Q}
}
\]
for a universal constant \(C\).

\subsection{The robust two-sample-size guarantee}
\label{sec:theory-main}

\begin{theorem}[Buffered ERM]
\label{thm:buffered-main}
Suppose Assumption~\ref{ass:comp-margin} holds and
\(2\varepsilon_N(\delta)\le t_0\).  Let
\(\widehat w_\varepsilon\) minimize
\(\widehat R_{\varepsilon_N}\) in Eq.~(\ref{eq:buffered-erm}).  With
probability at least \(1-\delta\) over both the \(Q\) questions and all
generations,
\begin{equation}
R(\widehat w_\varepsilon)-R(w^\star)
\le
2\mathfrak v_Q(\delta)
+
C_m(2\varepsilon_N(\delta))^\beta.
\label{eq:buffered-main-bound}
\end{equation}
Consequently,
\begin{equation}
R(\widehat w_\varepsilon)-R(w^\star)
\lesssim
\sqrt{\frac{(K-1)\log(eAQ)+\log(1/\delta)}{Q}}
+
C_m
\left(\frac{\log(KAQ/\delta)}{N}\right)^{\beta/2}.
\label{eq:buffered-rate}
\end{equation}
\end{theorem}

\begin{proof}
Write
\(\widetilde R_t(w)=Q^{-1}\sum_j\mathbf 1\{M_{q_j}(w)\le t\}\) and
\(R_t(w)=\mathbb P\{M_q(w)\le t\}\).  On the intersection of the events in
Lemma~\ref{lem:margin-accuracy} and Eq.~(\ref{eq:vq-event}),
\[
\begin{aligned}
R(\widehat w_\varepsilon)
&\le \widetilde R_0(\widehat w_\varepsilon)+\mathfrak v_Q \\
&\le \widehat R_{\varepsilon_N}(\widehat w_\varepsilon)+\mathfrak v_Q \\
&\le \widehat R_{\varepsilon_N}(w^\star)+\mathfrak v_Q \\
&\le \widetilde R_{2\varepsilon_N}(w^\star)+\mathfrak v_Q \\
&\le R_{2\varepsilon_N}(w^\star)+2\mathfrak v_Q.
\end{aligned}
\]
The remaining difference is exactly the positive comparator band:
\[
R_{2\varepsilon_N}(w^\star)-R(w^\star)
=
\mathbb P\{0<M_q(w^\star)\le2\varepsilon_N\}
\le C_m(2\varepsilon_N)^\beta.
\]
The two events have joint probability at least \(1-\delta\); substitution of
their definitions gives Eq.~(\ref{eq:buffered-rate}).
\end{proof}

The theorem separates two sources of uncertainty in learning the weights.  The
\(Q\)-term is ordinary generalization across questions.  The \(N\)-term is the
mass of comparator-correct training questions whose margins cannot be
distinguished from a tie at resolution \(N^{-1/2}\).  The latter can decay
faster than \(N^{-1/2}\) because the target is a thresholded correctness
decision rather than smooth parameter estimation.

\subsection{Why a margin condition does not imply a fast
\texorpdfstring{\(Q\)}{Q}-rate}
\label{sec:theory-obstruction}

The margin condition controls how much probability lies near a decision
boundary.  It does not say whether two well-separated regions of the weight
simplex have nearly identical risks.

\begin{proposition}[Hard margin with root-\(Q\) distinguishability]
\label{prop:counterexample}
There is a family with two models, two answers, \(N=\infty\), and a hard margin
at every optimum for which no learner can uniformly guarantee excess risk
\(o(Q^{-1/2})\) from \(Q\) labeled questions.
\end{proposition}

\begin{proof}
Write \(w=(t,1-t)\).  On type \(+\) questions, model 1 is always correct and
model 2 is always wrong, so \(M_q(w)=2t-1\).  On type \(-\) questions the
roles reverse and \(M_q(w)=1-2t\).  Consider distributions
\(\mathcal D_+\) and \(\mathcal D_-\) assigning type \(+\) probability
\(1/2+\eta/2\) and \(1/2-\eta/2\), respectively.  The optimal side of
\(t=1/2\) reverses between the two distributions; choosing the wrong side
costs excess risk \(\eta\).  At a pure optimal weight, every margin has
magnitude one, so there is no near-boundary mass.

Learning the correct side from \(Q\) questions is nevertheless testing two
coins with biases \(1/2\pm\eta/2\).  For
\(\eta=c/\sqrt Q\) with a sufficiently small constant \(c\), their
\(Q\)-sample laws have total variation bounded away from one.  Le Cam's
two-point argument therefore forces a constant probability of choosing the
wrong side under one of the distributions, and hence excess risk of order
\(Q^{-1/2}\).
\end{proof}

The obstruction is global identifiability, not boundary noise.  We therefore
state faster \(Q\)-rates only after adding a condition that risk grows when the
weights leave the optimum.

\subsection{Optional fast
\texorpdfstring{\(Q\)}{Q}-rates under identifiable geometry}
\label{sec:theory-fast}

\begin{assumption}[Absolute margin and risk growth]
\label{ass:growth}
There are \(C_m,\beta,t_0,c_g,\gamma,r_0,c_{\rm sep}>0\) such that
\begin{equation}
\mathbb P\{|M_q(w^\star)|\le t\}\le C_m t^\beta
\quad(0<t\le t_0),
\label{eq:absolute-margin}
\end{equation}
and
\begin{equation}
R(w)-R(w^\star)\ge
\begin{cases}
c_g\|w-w^\star\|_1^\gamma,
&\|w-w^\star\|_1\le r_0,\\
c_{\rm sep},
&\|w-w^\star\|_1>r_0.
\end{cases}
\label{eq:risk-growth}
\end{equation}
\end{assumption}

Let
\(D(w,w^\star)=\mathbb P\{\ell_q(w)\ne\ell_q(w^\star)\}\).

\begin{proposition}[Geometry implies a Bernstein condition]
\label{prop:geom-bernstein}
Under Assumption~\ref{ass:growth},
\begin{equation}
D(w,w^\star)
\le
B\{R(w)-R(w^\star)\}^{\alpha},
\qquad
\alpha=\min\{1,\beta/\gamma\},
\label{eq:bernstein-condition}
\end{equation}
for a finite constant \(B\).
\end{proposition}

\begin{proof}
Every contender gap is one-Lipschitz in \(\ell_1\), and taking a minimum
preserves that property.  If the losses at \(w\) and \(w^\star\) disagree,
then \(|M_q(w^\star)|\le\|w-w^\star\|_1\).  The absolute-margin condition
therefore gives
\(D(w,w^\star)\le C_m\|w-w^\star\|_1^\beta\) locally.  Inverting the local
risk-growth condition yields the exponent \(\beta/\gamma\); truncate it at
one because risks are bounded.  Outside the local ball, global separation and
\(D\le1\) complete the bound after enlarging the constant.  If
\(t_0<\|w-w^\star\|_1\le r_0\), the local growth condition itself gives
\(R(w)-R(w^\star)\ge c_gt_0^\gamma\), so this intermediate annulus is covered
by the same constant-enlargement argument.
\end{proof}

\begin{lemma}[Localized finite-growth bound]
\label{lem:localized-vc}
Let \(\mathcal L=\{\ell_q(w):w\in\Delta_K\}\) be a binary loss class with
minimizer \(w^\star\).  Suppose
\(\log\Pi_{\mathcal L}(n)\le V_n\) and \(V_{2Q}\le c_0V_Q\).  If
\[
\mathbb P\{\ell_q(w)\ne\ell_q(w^\star)\}
\le
B\{R(w)-R(w^\star)\}^{\alpha},
\qquad 0<\alpha\le1,
\]
and a data-dependent \(\widetilde w\) satisfies
\[
\frac1Q\sum_{j=1}^Q\ell_{q_j}(\widetilde w)
\le
\frac1Q\sum_{j=1}^Q\ell_{q_j}(w^\star)+a,
\]
then, with probability at least \(1-\delta\),
\begin{equation}
R(\widetilde w)-R(w^\star)
\le
C\left[
a+
\left(\frac{V_Q+\log(1/\delta)}{Q}\right)^{1/(2-\alpha)}
+\frac{V_Q+\log(1/\delta)}{Q}
\right],
\label{eq:localized-vc}
\end{equation}
where \(C\) depends only on \(B\), \(\alpha\), and \(c_0\).
\end{lemma}

\begin{proof}[Proof sketch]
For the excess loss \(g_w=\ell_q(w)-\ell_q(w^\star)\), the Bernstein premise
gives \(\mathbb E g_w^2\le B(\mathbb E g_w)^\alpha\).  Peeling the class into
shells \(r<\mathbb E g_w\le2r\), and applying Bernstein concentration with
the growth bound on each shell, gives deviation
\[
O\!\left(
\sqrt{\frac{r^\alpha\{V_Q+\log(1/\delta)\}}{Q}}
+\frac{V_Q+\log(1/\delta)}{Q}
\right).
\]
The fixed point in \(r\), together with the approximate-ERM slack \(a\),
yields Eq.~(\ref{eq:localized-vc}).  This is the standard relative
VC/localization argument \citep{massart2006risk,bartlett2005local,
koltchinskii2006local}; Appendix~\ref{app:fast-rate-proofs} gives the
specialization used here, including the \(2Q\)-sample growth step and the
data-dependent output.
\end{proof}

\begin{theorem}[Buffered ERM with an optional fast \(Q\)-rate]
\label{thm:fast}
Under Assumption~\ref{ass:growth}, let
\(\alpha=\min\{1,\beta/\gamma\}\).  For the same buffered estimator as in
Theorem~\ref{thm:buffered-main}, with probability at least \(1-\delta\),
\begin{equation}
R(\widehat w_\varepsilon)-R(w^\star)
\le
C
\left(
\frac{(K-1)\log(eAQ)+\log(1/\delta)}{Q}
\right)^{1/(2-\alpha)}
+
C'
\left(
\frac{\log(KAQ/\delta)}{N}
\right)^{\beta/2}.
\label{eq:fast-rate}
\end{equation}
The constants depend on the margin and growth parameters, not on
\(Q,N,K,\) or \(A\).
\end{theorem}

\begin{proof}[Proof sketch]
The indicator sandwich and optimality of the buffered estimator show that its
clean empirical risk exceeds that of \(w^\star\) by at most the empirical mass
of \(\{0<M_q(w^\star)\le2\varepsilon_N\}\).  Bernstein concentration for this
fixed band makes the slack
\(O(\varepsilon_N^\beta+\log(1/\delta)/Q)\).
Proposition~\ref{prop:geom-bernstein} supplies the Bernstein premise of
Lemma~\ref{lem:localized-vc}, with
\(V_Q=(K-1)\log(eAQ)\).  Applying that lemma to the preceding empirical slack
gives the localized \(Q\)-term; the band mass contributes the finite-\(N\)
term in Eq.~(\ref{eq:fast-rate}).  Appendix~\ref{app:fast-rate-proofs} tracks
the three confidence events and proves the full statement.
\end{proof}

The exponent is governed by two distinct geometric quantities.  \(\beta\)
measures the mass of questions that can change status near \(w^\star\);
\(\gamma\) measures how much the losses and gains fail to cancel when weights
move.  If crossings mostly hurt, then \(\gamma=\beta\), \(\alpha=1\), and the
\(Q\)-term is nearly \(Q^{-1}\).  If first-order losses and gains cancel, then
\(\gamma=\beta+1\), giving
\(\alpha=\beta/(\beta+1)\).  The latter exponent is a special geometry, not a
consequence of the margin condition by itself.
